% Test of keepafterheading and keepmaxgap (didactic): a heading keeps its lines,
% within the limit on empty space; where it can't, a warning.  In the notes
% (memoir + beamerarticle), with \raggedbottom, where penalties below 10000 don't
% count.
%
% Build: see the Makefile (LaTeX until the references settle).  The document
% checks itself: a line "didactic heading test: PASS/FAIL <check>" for each, then
% "didactic heading test: all checks ran"; the Makefile fails on any FAIL and on
% a missing last line, and counts the warnings.  Each case starts on a new page,
% after a rule that leaves room for a few lines (\roomfor).
%
%   AA  \section, \subsection, a sentence and a longtable that doesn't fit, with
%       more room than keepmaxgap: the headings stay together, the table goes to
%       the next page, and a warning says so.  (Within the limit, so that they
%       move: test-keep-nogap.)
%   AB  a frame title, a sentence and such a table.  AC \paragraph*, AD a block.
%   AE  a two-line heading, a one-line paragraph and a long one, room for all
%       of them to be moved: three lines after the heading on its page, when that
%       is what the option asks.
%   AF  the same with more room than keepmaxgap: it stays, and at keepafterheading=3
%       with a warning.  With a long paragraph, the same: it isn't moved past the limit.
%   AS  a listing after a listing: the lead-in of the first doesn't count in
%       what moving the second leaves empty, so it moves whole.
%   AG  first paragraphs with a footnote, a list, a float before them and a page
%       after them: only the first paragraph keeps its lines.
%   AH  a heading in a \loop: all iterations run, and scratch registers set
%       after a heading survive the start of its first paragraph.
%   AI  a heading, a lead-in and a listing, with more room than keepmaxgap when
%       counted from the heading: the listing breaks and the heading stays.
%   AJ  a lead-in and two listings in a row.  With room for the first but
%       for a few points, it breaks: nothing moves, and no page is overfull.
%       With room for the first, the second moves alone.
%   AK  a heading, a lead-in and a listing too short to break, with more room
%       than keepmaxgap at the heading: the tie to the lead-in is released, the
%       listing starts the next page, and a warning says so.
%   AL  a frame title, then a block or an example: keepafterheading lines after
%       the heading of the block.
%   AM  a heading, a one-line paragraph and a list: the estimate before the
%       heading moves them.  AN the same for a block.  AN2 with room for the block's
%       heading and that line, by the estimate: it stays and the list goes, warned.
%   AO  a paragraph at the foot of a page and at the top of the next: no last
%       line alone at the top; a first line may stay alone at the foot, the club
%       penalty being 1000.
%   AP  a table at the top of a page after a heading, with room to spare.
%   AR  a heading, a listing that takes no paragraph, then a paragraph: the
%       first paragraph after the listing isn't kept as a first paragraph.
%
% The default of keepafterheading here; test-keep-heading3 reads this file with
% keepafterheading=3, \testk.  The cases that depend on it are written with it.
\ifdefined\testk
  \PassOptionsToPackage{keepafterheading=\testk}{didactic}
\else
  \def\testk{2}
\fi
\documentclass[a4paper]{memoir}
\usepackage[british]{babel}
\usepackage{amsmath,amssymb,amsthm}
\usepackage{didactic}
\usepackage[noamsthm,notheorems]{beamerarticle}
\usepackage{refcount}
\usepackage{longtable}
\usepackage{lipsum}
\raggedbottom

\makeatletter
\newcommand\roomfor[1]{%
  \clearpage
  \noindent\rule{\linewidth}{0.4pt}\par
  \vspace*{\dimexpr\textheight-\topskip-#1\baselineskip-\baselineskip\relax}%
  \par}
\newcommand\testtable[1]{%
  \begin{longtable}{ll}
    \caption{A caption of several lines, several lines, several lines,
      several lines, several lines, several lines, several lines, several
      lines, several lines.}\label{#1}\\
    a & b\\ a & b\\ a & b\\ a & b\\
  \end{longtable}}
\newcommand\testresult[2]{%
  \typeout{didactic heading test: #2 #1}%
  \item #1: #2}
% \testcheck{name}{label1}{=|<}{label2}: same page, or the first on an earlier one.
\newcommand\testcheck[4]{%
  \ifnum\getpagerefnumber{#2}#3\getpagerefnumber{#4}\relax
    \testresult{#1}{PASS}%
  \else
    \testresult{#1 (#2 on p.\ \getpagerefnumber{#2}, #4 on p.\ \getpagerefnumber{#4})}{FAIL}%
  \fi}
% \testclubs{name}{n}: the number of \clubpenalties here, while the check list is
% collected at the end.
\newcommand\testlist{}
\newcommand\testclubs[2]{%
  \edef\@tempa{\the\clubpenalties 0}%
  \ifnum\@tempa=#2\relax
    \g@addto@macro\testlist{\testresult{#1}{PASS}}%
  \else
    \xdef\@gtempa{#1, n=\@tempa, not #2}%
    \expandafter\testfail\expandafter{\@gtempa}%
  \fi}
\newcommand\testfail[1]{\g@addto@macro\testlist{\testresult{#1}{FAIL}}}
% \testtrue{name}{condition}
\newcommand\testif[2]{%
  #2%
    \g@addto@macro\testlist{\testresult{#1}{PASS}}%
  \else
    \testfail{#1}%
  \fi}
% \testlisting{prefix}{n}: a listing of n lines, each with a label.
\newwrite\@test@out
\newcount\@test@i
\newcommand\testlisting[2]{%
  \immediate\openout\@test@out=\jobname-#1.txt
  \@test@i\z@
  \loop
    \advance\@test@i\@ne
    \immediate\write\@test@out{line \the\@test@i\string\label{#1L\the\@test@i}}%
  \ifnum\@test@i<#2\repeat
  \immediate\closeout\@test@out
  \VerbatimInput[commandchars=\\\{\}]{\jobname-#1.txt}}
\makeatother
\newcommand\oneline{One line.}
\newcommand\nl{\hspace*{\fill}\break}
\newcommand\threelines{First paragraph, with at least three lines of text in it, and on it
  goes, and on it goes, and on it goes, and on it goes, and on it goes, and on it
  goes, and on it goes, and on it goes, and on it goes, and on it goes.}
\newcommand\longpara{Second paragraph, long enough to break across pages. \lipsum[2]}

\begin{document}

\roomfor{5}
\section*{AA section}\label{aa-a}
\subsection*{AA subsection}\label{aa-b}
Table lists the rows.
\testtable{aa-t}
Text after AA.

\mode<article>{\roomfor{4}}
\begin{frame}{AB frame title}
  \label{ab-a}Table lists the rows.
  \testtable{ab-t}
\end{frame}

\roomfor{4}
\paragraph*{AC paragraph}\label{ac-a}
Table lists the rows.
\testtable{ac-t}
Text after AC.

\roomfor{4}
\begin{example}[AD]
  \label{ad-a}Table lists the rows.
  \testtable{ad-t}
\end{example}

\roomfor{3}
\section*{AE, a heading long enough to run over two lines of its large font,
  indeed it is}\label{ae-a}
\oneline

First line\label{ae-1}\\
second line\label{ae-2}\\
third line\\
fourth line.

\roomfor{6}
\section*{AF, a heading long enough to run over two lines of its large font,
  indeed it is}\label{af-a}
\oneline

\longpara

\roomfor{6}\label{ag-r}
\section*{AF, a heading long enough to run over two lines of its large font,
  indeed it is}\label{ag-a}
First line of a long paragraph\nl
second line\nl
third line\nl
fourth line\nl
fifth line.

\roomfor{8.5}
Lead\\in\\three.
\testlisting{u}{1}
\testlisting{v}{10}
Text after AS.

\roomfor{30}
\section*{AG, a footnote}
\threelines\footnote{A note.}\testclubs{AG first paragraph with a footnote}{\testk}

Second.\testclubs{AG second paragraph after a footnote}{0}

Third.\testclubs{AG third paragraph after a footnote}{0}

\section*{AG, a list}
\threelines\testclubs{AG first paragraph before a list}{\testk}
\begin{itemize}\item Item.\testclubs{AG item after the first paragraph}{0}\end{itemize}
Second.\testclubs{AG paragraph after a list}{0}

\section*{AG, a float}
\begin{figure}[t]\centering FIGURE\caption{A figure.}\end{figure}
\threelines\testclubs{AG first paragraph after a float}{\testk}

Second.\testclubs{AG second paragraph after a float}{0}

\section*{AG, a page}
\threelines\testclubs{AG first paragraph before a page}{\testk}

\clearpage
Second.\testclubs{AG second paragraph after a page}{0}

Third.\testclubs{AG third paragraph after a page}{0}

\section*{AG, a list first}
\begin{itemize}
\item \threelines\testclubs{AG a list as the first thing}{\testk}
\item Second item.\testclubs{AG the second item}{0}
\end{itemize}
Text.\testclubs{AG paragraph after that list}{0}

\clearpage
\newcount\aliter
\aliter=0
\loop
  \advance\aliter by 1
  \section*{AH \the\aliter}
  Iteration \the\aliter. \threelines\par
\ifnum\aliter<3 \repeat
\makeatletter
\testif{AH loop ran three times}{\ifnum\aliter=3 }
\section*{AH scratch}
\count@=42 \dimen@=42pt \def\reserved@a{kept}%
\threelines\par
\testif{AH count@ survives}{\ifnum\count@=42 }
\testif{AH dimen@ survives}{\ifdim\dimen@=42pt }
\def\@tempa{kept}%
\testif{AH reserved@a survives}{\ifx\reserved@a\@tempa}
\makeatother

\roomfor{7}
\section*{AI}\label{ai-a}
Lead-in AI.\label{ai-l}
\testlisting{i}{10}
Text after AI.

\roomfor{10}
Lead-in AJ.\label{aj-l}
\testlisting{j}{10}
\testlisting{k}{8}
Text after AJ.

\roomfor{12}
Lead-in AJ.\label{aj2-l}
\testlisting{s}{10}
\testlisting{t}{8}
Text after AJ.

\roomfor{7}
\section*{AK}\label{ak-a}
Lead-in AK.\label{ak-l}
\testlisting{m}{4}
Text after AK.

\mode<article>{\roomfor{4}}
\begin{frame}{AL frame title}
  \begin{block}{AL block}
    First line\label{al-1}\nl
    second line\label{al-2}\nl
    third line\label{al-3}\nl
    fourth line.
  \end{block}
\end{frame}

\mode<article>{\roomfor{4}}
\begin{frame}{AL frame title}
  \begin{example}[AL]
    First line\label{am-1}\nl
    second line\label{am-2}\nl
    third line\label{am-3}\nl
    fourth line.
  \end{example}
\end{frame}

\roomfor{3.7}
\section*{AM}\label{an-a}
\oneline
\begin{itemize}
\item\rule{0pt}{5\baselineskip}A tall item.\label{an-i}
\item Another item.
\end{itemize}

\roomfor{\ifnum\testk=2 1.2\else 2.2\fi}
\begin{example}[AN]\label{ao-a}
  \oneline
  \begin{itemize}
  \item\rule{0pt}{5\baselineskip}A tall item.\label{ao-i}
  \item Another item.
  \end{itemize}
\end{example}

\roomfor{\ifnum\testk=2 1.6\else 2.6\fi}\label{ao2-r}
\begin{example}[AN2]\label{ao2-a}
  \oneline
  \begin{itemize}
  \item\rule{0pt}{5\baselineskip}A tall item.\label{ao2-i}
  \item Another item.
  \end{itemize}
\end{example}

\roomfor{0}
First line\label{ap-1}\\
second line\label{ap-2}\\
third line\\
fourth line\\
fifth line.

\roomfor{3}
First line\\
second line\\
third line\\
fourth line\label{aq-4}\\
fifth line.\label{aq-5}

\clearpage
\section*{AP}\label{ar-a}
Table lists the rows.
\testtable{ar-t}

\clearpage
\section*{AR}
\testlisting{r}{2}
\testclubs{AR listing took the state}{0}
Paragraph after the listing.\testclubs{AR paragraph after a listing}{0}

\clearpage
\section*{Checks}
\begin{itemize}
\testcheck{AA outer heading with the inner one}{aa-a}{=}{aa-b}
\testcheck{AA table on the next page}{aa-b}{<}{aa-t}
\testcheck{AB frame title with the table}{ab-a}{=}{ab-t}
\testcheck{AC paragraph with the table}{ac-a}{=}{ac-t}
\testcheck{AD block with the table}{ad-a}{=}{ad-t}
\testcheck{AE heading with its second line}{ae-a}{=}{ae-1}
\ifnum\testk>2
\testcheck{AE heading with its third line}{ae-a}{=}{ae-2}
\fi
\testcheck{AF long paragraph: heading stays}{ag-r}{=}{ag-a}
\testcheck{AS second listing whole}{vL1}{=}{vL10}
\testcheck{AI heading with the listing's first line}{ai-a}{=}{iL1}
\testcheck{AI lead-in with the first three lines}{ai-l}{=}{iL3}
\testcheck{AI listing breaks}{iL1}{<}{iL10}
\testcheck{AJ lead-in with the first three lines}{aj-l}{=}{jL3}
\testcheck{AJ first listing breaks}{jL3}{<}{jL10}
\testcheck{AJ first listing with its lead-in, with room}{aj2-l}{=}{sL10}
\testcheck{AJ second listing on the next page}{sL10}{<}{tL1}
\testcheck{AJ second listing whole}{tL1}{=}{tL8}
\testcheck{AK heading stays with its lead-in}{ak-a}{=}{ak-l}
\testcheck{AK listing on the next page}{ak-l}{<}{mL1}
\testcheck{AK listing whole}{mL1}{=}{mL4}
\testcheck{AL block's lines (block)}{al-1}{=}{al-\testk}
\testcheck{AL block's lines (example)}{am-1}{=}{am-\testk}
\testcheck{AM heading with the list}{an-a}{=}{an-i}
\testcheck{AN block with the list}{ao-a}{=}{ao-i}
\testcheck{AN2 block stays with its line}{ao2-r}{=}{ao2-a}
\testcheck{AN2 the list on the next page}{ao2-a}{<}{ao2-i}
\testcheck{AO first line alone at the foot}{ap-1}{<}{ap-2}
\testcheck{AO last line not alone}{aq-4}{=}{aq-5}
\testcheck{AP heading with the table}{ar-a}{=}{ar-t}
\testlist
\end{itemize}
\typeout{didactic heading test: all checks ran}
\end{document}
